\chapter{The King and the Lieutenant}

\lettrine{A}{s} soon as the king saw the officer enter, he dismissed his valet de chambre and his gentleman.

\zz
<Who is on duty to-morrow, monsieur?> asked he.

\zz
The lieutenant bowed his head with military politeness, and replied, <I am, sire.>

<What! still you?>

<Always I, sire.>

<How can that be, monsieur?>

<Sire, when travelling, the musketeers supply all the posts of your majesty's household; that is to say, yours, her majesty the queen's, and monsieur le cardinal's, the latter of whom borrows of the king the best part, or rather the numerous part, of the royal guard.>

<But in the interims?>

<There are no interims, sire, but for twenty or thirty men who rest out of a hundred and twenty. At the Louvre it is very different, and if I were at the Louvre I should rely upon my brigadier; but, when travelling, sire, no one knows what may happen, and I prefer doing my duty myself.>

<Then you are on guard every day?>

<And every night. Yes, sire.>

<Monsieur, I cannot allow that—I will have you rest.>

<That is very kind, sire; but I will not.>

<What do you say?> said the king, who did not at first comprehend the full meaning of this reply.

<I say, sire, that I will not expose myself to the chance of a fault. If the devil had a trick to play on me, you understand, sire, as he knows the man with whom he has to deal, he would chose the moment when I should not be there. My duty and the peace of my conscience before everything, sire.>

<But such duty will kill you, monsieur.>

<Eh! sire, I have performed it for thirty years, and in all France and Navarre there is not a man in better health than I am. Moreover, I entreat you, sire, not to trouble yourself about me. That would appear very strange to me, seeing that I am not accustomed to it.>

The king cut short the conversation by a fresh question. <Shall you be here, then, to-morrow morning?>

<As at present? yes, sire.>

The king walked several times up and down his chamber; it was very plain that he burned with a desire to speak, but that he was restrained by some fear or other. The lieutenant, standing motionless, hat in hand, watched him making these evolutions, and, whilst looking at him, grumbled to himself, biting his moustache:

<He has not half a crown worth of resolution! Parole d'honneur! I would lay a wager he does not speak at all!>

The king continued to walk about, casting from time to time a side glance at the lieutenant. <He is the very image of his father,> continued the latter, in is secret soliloquy, <he is at once proud, avaricious, and timid. The devil take his master, say I\@.>

The king stopped. <Lieutenant,> said he.

<I am here, sire.>

<Why did you cry out this evening, down below in the salons—<The king's service! His majesty's musketeers!>>

<Because you gave me the order, sire.>

<I\@?> <Yourself.>

<Indeed, I did not say a word, monsieur.>

<Sire, an order is given by a sign, by a gesture, by a glance, as intelligibly, as freely, and as clearly as by word of mouth. A servant who has nothing but ears is not half a good servant.>

<Your eyes are very penetrating, then, monsieur.>

<How is that, sire?>

<Because they see what is not.>

<My eyes are good, though, sire, although they have served their master long and much: when they have anything to see, they seldom miss the opportunity. Now, this evening, they saw that your majesty coloured with endeavouring to conceal the inclination to yawn, that your majesty looked with eloquent supplications, first to his eminence, and then at her majesty, the queen-mother, and at length to the entrance door, and they so thoroughly remarked all I have said, that they saw your majesty's lips articulate these words: <Who will get me out of this?>>

<Monsieur!>

<Or something to this effect, sire—<My musketeers!> I could then no longer hesitate. That look was for me. I cried out instantly, <His majesty's musketeers!> And, besides, that was shown to be true, sire, not only by your majesty's not saying I was wrong, but proving I was right by going out at once.>

The king turned away to smile; then, after a few seconds, he again fixed his limpid eye upon that countenance, so intelligent, so bold, and so firm, that it might have been said to be the proud and energetic profile of the eagle facing the sun. <That is all very well,> said he, after a short silence, during which he endeavoured, in vain, to make his officer lower his eyes.

But seeing the king said no more, the latter pirouetted on his heels, and took three steps towards the door, muttering, <He will not speak! \swear{Mordioux!} he will not speak!>

<Thank you, monsieur,> said the king at last.

<Humph!> continued the lieutenant; <there was only wanting that. Blamed for having been less of a fool than another might have been.> And he went to the door, allowing his spurs to jingle in true military style. But when he was on the threshold, feeling the king's desire drew him back, he returned.

<Has your majesty told me all?> asked he, in a tone we cannot describe, but which, without appearing to solicit the royal confidence, contained so much persuasive frankness, that the king immediately replied:

<Yes; but draw near, monsieur.>

<Now then,> murmured the officer, <he is coming to it at last.>

<Listen to me.>

<I shall not lose a word, sire.>

<You will mount on horseback to-morrow, at about half-past four in the morning, and you will have a horse saddled for me.>

<From your majesty's stables?>

<No; one of your musketeers' horses.>

<Very well, sire. Is that all?>

<And you will accompany me.>

<Alone?>

<Alone.>

<Shall I come to seek your majesty, or shall I wait?>

<You will wait for me.>

<Where, sire?>

<At the little park-gate.>

The lieutenant bowed, understanding that the king had told him all he had to say. In fact, the king dismissed him with a gracious wave of the hand. The officer left the chamber of the king, and returned to place himself philosophically in his fauteuil, where, far from sleeping, as might have been expected, considering how late it was, he began to reflect more deeply than he had ever reflected before. The result of these reflections was not so melancholy as the preceding ones had been.

<Come, he has begun,> said he. <Love urges him on, and he goes forward—he goes forward! The king is nobody in his own palace; but the man perhaps may prove to be worth something. Well, we shall see to-morrow morning. Oh! oh!> cried he, all at once starting up, <that is a gigantic idea, \swear{mordioux!} and perhaps my fortune depends, at least, upon that idea!> After this exclamation, the officer arose and marched, with his hands in the pockets of his \binderyfrench{justaucorps}, about the immense ante-chamber that served him as an apartment. The wax-light flamed furiously under the effects of a fresh breeze, which stole in through the chinks of the door and the window, and cut the salle diagonally. It threw out a reddish, unequal light, sometimes brilliant, sometimes dull, and the tall shadow of the lieutenant was seen marching on the wall, in profile, like a figure by Callot, with his long sword and feathered hat.

<Certainly!> said he, <I am mistaken if Mazarin is not laying a snare for this amorous boy. Mazarin, this evening, gave an address, and made an appointment as complacently as M.~Dangeau himself could have done—I heard him, and I know the meaning of his words. <To-morrow morning,> said he, <they will pass opposite the bridge of Blois.> \swear{Mordioux!} that is clear enough, and particularly for a lover. That is the cause of this embarrassment; that is the cause of this hesitation; that is the cause of this order—<Monsieur the lieutenant of my musketeers, be on horseback to-morrow at four o'clock in the morning.> Which is as clear as if he had said,—<Monsieur the lieutenant of my musketeers, to-morrow, at four, at the bridge of Blois,—do you understand?> Here is a state secret, then, which I, humble as I am, have in my possession, while it is in action. And how do I get it? Because I have good eyes, as his majesty just now said. They say he loves this little Italian doll furiously. They say he threw himself at his mother's feet, to beg her to allow him to marry her. They say the queen went so far as to consult the court of Rome, whether such a marriage, contracted against her will, would be valid. Oh, if I were but twenty-five! If I had by my side those I no longer have! If I did not despise the whole world most profoundly, I would embroil Mazarin with the queen-mother, France with Spain, and I would make a queen after my own fashion. But let that pass.> And the lieutenant snapped his fingers in disdain.

<This miserable Italian—this poor creature—this sordid wretch—who has just refused the king of England a million, would not perhaps give me a thousand pistoles for the news I would carry him. \swear{Mordioux!} I am falling into second childhood—I am becoming stupid indeed! The idea of Mazarin giving anything! ha! ha! ha!> and he laughed in a subdued voice.

<Well, let us go to sleep—let us go to sleep; and the sooner the better. My mind is wearied with my evening's work, and will see things to-morrow more clearly than to-day.>

And upon this recommendation, made to himself, he folded his cloak around him, looking with contempt upon his royal neighbour. Five minutes after this he was asleep, with his hands clenched and his lips apart, giving escape, not to his secret, but to a sonorous sound, which rose and spread freely beneath the majestic roof of the ante-chamber. 